\subsection{Observed requests in real agent runs}
\label{sec:agent-observations}

\paragraph{Capture and scope.}
We reanalyzed eight SWE-bench Verified task runs~\cite{swebench2024}: five
Core and three Exploratory. All used the same agent harness configured for
\code{deepseek-v4-flash}. Exploratory runs used a broader
investigation/testing prompt and a 30-minute budget; all finished earlier.
Each base-commit checkout was mounted through an instrumented passthrough
file system on Ubuntu 24.04 under WSL, recording actual FUSE handler
requests. The companion record retains the harness and instrumentation versions.
Streaming reanalysis covered all 1,138,884 events: 217,129 Core and 921,755
Exploratory. Counts retain every status, 64 setup/canary requests before
agent startup, and teardown, including 24 requests after recorded run finish.
Group statistics weight requests equally, not tasks. Per-task counts and
source-validation details are in the companion workload record.

\begin{figure}[!t]
\centering
\includegraphics[width=\columnwidth]{figures/agent-fuse/fuse-single-column.pdf}
\caption{FUSE request structure across eight task runs.
(a) All-call shares; hatched segments are errors.
(b) Requested-length distributions within each method and group.
(c) Each task's call share in its busiest 10\% of one-second windows.
Core precedes Exploratory in each bar pair.}
\label{fig:agent-fuse-combined}
\end{figure}

\paragraph{Metadata traffic and read dominance.}
The metadata methods---\code{lookup}, \code{getattr}, \code{getxattr},
\code{readdirplus}, \code{opendir}, and \code{releasedir}---accounted for
61.5\% of Core and 56.4\%
of Exploratory calls (Figure~\ref{fig:agent-fuse-combined}a).
Other comprises all methods except these six and read/write.
Restricting the denominator and numerator to \code{ok} events reduces these
shares to 52.0\% and 44.2\%: failed probes materially change the request mix.
Across tasks, the all-status share ranged from 39.6\% in SymPy to 65.4\% in
Django, so the group majority does not describe every task.
Read calls outnumbered writes by 19.9:1 in Core and 4.16:1 in Exploratory;
per-task ratios ranged from 3.12:1 (Pytest) to 42.7:1 (Matplotlib).
These observations motivate evaluating both positive and negative metadata
paths alongside data I/O.

\paragraph{A mixture of request sizes.}
All 22,142/98,292 reads and 1,112/23,600 writes (Core/Exploratory) were
\code{ok}; $1\,\mathrm{KiB}=1024$ bytes.
Reads ranged from 4 to 128 KiB, with Core and Exploratory medians of 8 and
16 KiB. Writes ranged from one byte
to 128 KiB, with medians of 4,084 and 4,117 bytes, respectively.
Requests of at most 4 KiB accounted for 47.7\% and 34.5\% of reads, and
74.3\% and 49.9\% of writes. The interval $(4,64]$ KiB contained a further
36.2\% and 36.4\% of reads and 21.3\% and 44.0\% of writes
(Figure~\ref{fig:agent-fuse-combined}b).
This mixture motivates measuring how client fetching and uploads transform
requests before testing backend layouts. FUSE observations occur below kernel
caching and readahead; they do not identify application syscall sizes, unique
bytes, finalized object lengths, or an optimal block size.

\paragraph{Task and temporal heterogeneity.}
Tasks generated 8,542--634,726 calls. Astropy contributed 68.9\% of
Exploratory calls, versus 6.0\% for Pytest. Capture spans, from first to last
request start, were 42.7--71.3 seconds for Core and 119.9--556.1 seconds for
Exploratory. Windows include zero-call and final partial windows.
The busiest 10\% of one-second windows, rounded up per task,
contained 47.4--83.7\% of its calls (Figure~\ref{fig:agent-fuse-combined}c).
The companion record retains the full method breakdowns, size ECDFs,
and per-task cumulative time curves. Replay
should retain these per-task mixes and bursts. Each task was observed once;
different repositories, prompts, and budgets prevent a causal interpretation
of the group contrast. Without tool-phase labels, bursts cannot be assigned
to searching, editing, or testing.

\paragraph{Status semantics and evidence limits.}
The capture contained 872,519 \code{ok}, 253,669 \code{error}, and 12,696
\code{incomplete} events. Errors occurred in \code{lookup} (233,452),
\code{getxattr} (20,203), and \code{rmdir} (14). Collection notes associate
many with failed path probes or unsupported attributes; absent errno values
prevent event-level cause attribution. The parser labels a request
\code{incomplete} when no recognized completion is paired with it: here,
12,688 \code{forget} and eight \code{destroy} records. Neither label counts
failed agent tasks. All eight agent and daemon exit codes were zero, which
does not establish SWE-bench grading success.

Handler latency measures WSL passthrough intervals and excludes model response
time. Elapsed capture time instead includes inter-request gaps, which may
contain model waiting; summing overlapping handler intervals cannot recover
task wall time. These observations do not measure RustFS, Redis, network
object storage, or cloud-disk latency. They motivate metadata and mixed-I/O
experiments for this single-agent/model panel, but do not establish cross-task
sharing, sparse mutation, or performance gains from \sysname, RustFS, or
the automatic hint policy. Those remain questions for
Section~\ref{sec:evaluation}.
